\section{Entidades y relaciones}

\subsection{Descripcion}

Este documento enumera las entidades conceptuales del sistema y sus relaciones principales. El nivel es funcional y no pretende reemplazar un esquema fisico de base de datos.

\subsection{Alcance}

Aplica al modelo conceptual de datos de la plataforma.

\subsection{Definiciones}

\begin{itemize}
  \item \textbf{Entidad}: objeto de dominio con identidad propia.
  \item \textbf{Relacion}: vinculo funcional entre entidades.
\end{itemize}

\subsection{Reglas}

Entidades principales:

\begin{longtable}{p{0.25\textwidth}p{0.67\textwidth}}
\toprule
\textbf{Entidad} & \textbf{Responsabilidad funcional} \\
\midrule
User & Identidad del participante o administrador. \\
PaymentProof & Evidencia de pago usada para activacion manual. \\
Tournament & Contenedor del torneo documentado por la plataforma. \\
Phase & Segmentacion oficial del torneo para scoring y vistas. \\
Team & Equipo participante del torneo. \\
Match & Partido con kickoff, fase y estado operativo. \\
Pick & Prediccion de marcador exacto por usuario y partido. \\
SpecialPrediction & Prediccion especial de Campeon o Goleador dentro del subdominio de picks. \\
Result & Resultado oficial o corregido de un partido. \\
ScoreEntry & Registro derivado del puntaje aplicado a un pick. \\
LeaderboardSnapshot & Vista derivada opcional del ranking en un momento dado. \\
AuditEvent & Evento auditable generado por accion critica o correccion. \\
\bottomrule
\end{longtable}

Notas de modelado conceptual:

\begin{itemize}
  \item `Pick` y `SpecialPrediction` se documentan como conceptos distintos porque tienen distinta cardinalidad, distinta ventana de lock y distinta forma de evaluacion.
  \item Esta distincion conceptual no obliga por si misma a una implementacion fisica especifica en base de datos.
  \item `LeaderboardSnapshot` se mantiene como entidad derivada opcional; la fuente normativa del ranking sigue siendo el scoring reproducible.
\end{itemize}

Relaciones principales:

\begin{itemize}
  \item `User` genera `PaymentProof`.
  \item `User` crea `Pick` y `SpecialPrediction`.
  \item `Tournament` contiene `Phase`, `Team` y `Match`.
  \item `Tournament` define tambien el contexto unico de cada `SpecialPrediction`.
  \item `Match` tiene cero o un `Result` vigente.
  \item `Pick` se evalua contra `Result` para producir `ScoreEntry`.
  \item `SpecialPrediction` se evalua contra el cierre oficial del torneo o contra su estadistica oficial correspondiente para producir puntaje acumulable.
  \item `AuditEvent` puede referenciar acciones sobre usuarios, pagos, partidos, resultados y recalculos.
\end{itemize}

\subsection{Casos borde}

\begin{itemize}
  \item Un `Match` puede pasar por estados no lineales como postergado o corregido; el modelo debe preservar historia suficiente para auditar.
  \item Un `User` puede permanecer en \texttt{PENDIENTE} mientras su `PaymentProof` atraviesa revision administrativa sin que eso implique nuevos estados funcionales del usuario.
  \item `LeaderboardSnapshot` es opcional; el ranking puede derivarse o persistirse segun necesidad operativa.
\end{itemize}

\subsection{Invariantes}

\begin{itemize}
  \item Un `Pick` pertenece a un solo `User` y a un solo `Match`.
  \item Una `SpecialPrediction` pertenece a un solo `User`, a un solo `Tournament` y a un solo tipo de prediccion.
  \item Un `Result` vigente a efectos de scoring debe ser unico por `Match`.
  \item Cada `AuditEvent` debe poder asociarse a un actor o a un proceso identificable.
\end{itemize}

\subsection{Preguntas abiertas}

\begin{itemize}
  \item Confirmar si `ScoreEntry` registrara tambien de forma explicita el puntaje de `SpecialPrediction` o si ese puntaje vivira en una estructura derivada distinta.
\end{itemize}
